\scanpage{24}
(\unsure{c|b}) \deleted{\unsure{Soit}} Le noyau de $\varphi\mid D_n^0$ est un sous gpe de $D_n^0=Z/n\mathbb Z$, donc \deleted{de la} c'est le ss-groupe engendré par un $u^i$ où $i$ est un diviseur \unsure{unqt|un\textsuperscript{t}|$>0$} déterminé de $n$ -- \deleted{\unsure{c|et}} qui a priori pourrait être $1$ ou $2$. Mais si \deleted{\unsure{$n=$|$i$}} $n'$ est tel que
\[
n=n'd
\]
\deleted{$\varphi$ se factorise en}
\[
\deleted{\varphi':D_{n'}\mathbin{\unsure{-|\to}}\mathrm{Aff}(E)}
\]
on voit que \deleted{$S$ a un} $\operatorname{card}S=n'$, \deleted{donc} et comme $S$ engendre aff\textsuperscript{t} $E$, on aura $n'\geq3$. Enfin, comme $\varphi$ se factorise par
\[
\varphi':D_{n'}\mathbin{\unsure{-|\to}}\mathrm{Aff}(E)
\]
avec $\varphi'\mid D_n^0$ \unsure{$^{\bullet}$} fidèle -- et que les conditions (\unsure{1|11}) sont évidemment satisfaites aussi pour $\unsure{\rho'|\varphi'}$ -- (b) implique que $\varphi'$ est fidèle, donc $\operatorname{Ker}\varphi^{\unsure{\prime\prime|\circ}}$ est ident à $\operatorname{Ker}(\varphi\mid D_n^0)$.

\margin{se montre avant (c)}
(b) \deleted{\unsure{Résulte du|Résulte de}} \deleted{Supposons $\varphi\mid D_n^0$ fidèle.} Équivaut au
% archive sheet number: 23
